\chapter{Equazioni diofantee e congruenziali}\label{cap:diofantee}

Con il teorema di Bézout e gli anelli $\mbZ_m$ possiamo risolvere due tipi di equazioni a coefficienti interi: le
\textbf{equazioni diofantee} $ax + by = c$, in cui si cercano soluzioni intere, e le \textbf{equazioni congruenziali} $ax \equiv_m b$. Vedremo
che sono due facce dello stesso problema, e ne ricaveremo chi sono gli elementi invertibili di $\mbZ_m$. Chiudiamo con il
\textbf{periodo} di un elemento di un gruppo, che dice quando due potenze dello stesso elemento coincidono.

% ──────────────────────────────────────────────────────────────
\section{Equazioni diofantee}\label{sec:diofantee}

\dfn[diofantea]{Equazione diofantea di primo grado}{
    Siano $a, b, c \in \mbZ$. La funzione $(x, y) \in \mbZ \times \mbZ \mapsto ax + by - c \in \mbZ$ si chiama \textbf{equazione diofantea} di primo
    grado in due incognite, di termini $a, b, c$; per brevità si scrive $ax + by = c$. Una coppia $(m, n) \in \mbZ \times \mbZ$ si dice
    \textbf{soluzione} se $am + bn = c$.
}

Le combinazioni di $a$ e $b$ sono proprio gli elementi del sottogruppo che generano.

\begin{concetto}
\begin{Lemma}{Il sottogruppo generato da due interi}{sottogruppo-ab}
    \begin{prerequisiti}
        \dip{Definizione}{def:generata}
        \dip{Definizione}{def:sottostruttura}
    \end{prerequisiti}
    Siano $a, b \in \mbZ$. Nel gruppo $(\mbZ, +)$ il sottogruppo generato da $\{a, b\}$ è $\langle a, b \rangle = \{au + bv \mid u, v \in \mbZ\}$.
\end{Lemma}
\begin{dimostrazione}
    L'insieme $c = \{au + bv \mid u, v \in \mbZ\}$ è un sottogruppo: contiene $0 = a0 + b0$, la somma di due combinazioni è una combinazione
    e l'opposto di $au + bv$ è $a(-u) + b(-v)$. Contiene $a = a1 + b0$ e $b = a0 + b1$, quindi $\langle a, b \rangle \subseteq c$. Viceversa ogni
    sottogruppo che contiene $a$ e $b$ contiene i loro multipli e le loro somme, cioè tutte le combinazioni: $c \subseteq \langle a, b \rangle$.
\end{dimostrazione}
\end{concetto}

\begin{spiegazione}
    Con monete da $a$ e da $b$, potendo anche «restituire» monete (coefficienti negativi), si ottengono tutte le somme $au + bv$.
    Ad esempio con $4$ e $6$: $6 - 4 = 2$, e quindi tutti i numeri pari, e nient'altro: $\langle 4, 6 \rangle = 2\mbZ$.
\end{spiegazione}

\begin{concetto}
\begin{Teorema}{Forme equivalenti del teorema di Bézout}{bezout-equivalenze}
    \begin{prerequisiti}
        \dip{Teorema}{teo:bezout}
        \dip{Lemma}{lem:sottogruppo-ab}
        \dip{Definizione}{def:coprimi}
        \dip{Proposizione}{prop:prop-divisibilita}
        \dip{Definizione}{def:diofantea}
    \end{prerequisiti}
    Siano $a, b \in \mbZ \setminus \{0\}$ e $d \in \mathrm{MCD}(a, b)$. Gli enunciati seguenti (letti per tutte le coppie $a, b$) sono equivalenti, e
    quindi tutti veri:
    \begin{enumerate}
        \item (\textbf{Bézout}) esistono $u, v \in \mbZ$ con $d = au + bv$;
        \item $a$ e $b$ sono coprimi se e solo se esistono $u, v \in \mbZ$ con $1 = au + bv$;
        \item in $(\mbZ, +)$ si ha $\langle a, b \rangle = d\mbZ$;
        \item l'equazione diofantea $ax + by = c$ ha soluzioni se e solo se $d \mid c$.
    \end{enumerate}
\end{Teorema}
\begin{dimostrazione}
    \begin{caso}{(1) $\rightarrow$ (2)}
        $(\rightarrow)$ Se $a, b$ sono coprimi, $1 \in \mathrm{MCD}(a, b)$ e si applica (1) con $d = 1$. $(\leftarrow)$ Se $1 = au + bv$, ogni divisore comune
        di $a$ e $b$ divide $1$ (Proposizione~\ref{prop:prop-divisibilita}), quindi è $\pm 1$: $a$ e $b$ sono coprimi.
    \end{caso}
    \begin{caso}{(2) $\rightarrow$ (3)}
        $(\subseteq)$ $a, b \in d\mbZ$, che è un sottogruppo, quindi $\langle a, b \rangle \subseteq d\mbZ$. $(\supseteq)$ Scriviamo $a = a_1 d$, $b = b_1 d$.
        $a_1$ e $b_1$ sono coprimi: se $e$ dividesse entrambi con $e \neq \pm 1$, $ed$ sarebbe un divisore comune di $a$ e $b$ che non divide $d$. Per (2)
        $1 = a_1 u + b_1 v$, e moltiplicando per $d$: $d = au + bv \in \langle a, b \rangle$. Quindi $d\mbZ \subseteq \langle a, b \rangle$.
    \end{caso}
    \begin{caso}{(3) $\rightarrow$ (4)}
        L'equazione ha soluzioni se e solo se $c$ è una combinazione di $a$ e $b$, cioè (Lemma~\ref{lem:sottogruppo-ab}) $c \in \langle a, b \rangle = d\mbZ$,
        cioè $d \mid c$.
    \end{caso}
    \begin{caso}{(4) $\rightarrow$ (1)}
        Con $c = d$: poiché $d \mid d$, l'equazione $ax + by = d$ ha una soluzione $(u, v)$, cioè $d = au + bv$.
    \end{caso}
\end{dimostrazione}
\end{concetto}

\begin{spiegazione}
    I quattro enunciati sono quattro modi di dire la stessa cosa: \textbf{le combinazioni intere di $a$ e $b$ sono esattamente i multipli
    del loro MCD}. Per $a = 4$ e $b = 6$, con $d = 2$:
    \begin{passi}
        \item Bézout: $2 = 4 \cdot (-1) + 6 \cdot 1$;
        \item $4$ e $6$ non sono coprimi, e infatti $1$ non è una combinazione di $4$ e $6$ (le combinazioni sono tutte pari);
        \item $\langle 4, 6 \rangle = 2\mbZ$;
        \item $4x + 6y = 10$ ha soluzioni ($2 \mid 10$: ad esempio $x = 1$, $y = 1$), mentre $4x + 6y = 5$ no ($2 \nmid 5$).
    \end{passi}
\end{spiegazione}

\begin{concetto}
\begin{Teorema}{Tutte le soluzioni di un'equazione diofantea}{soluzioni-diofantee}
    \begin{prerequisiti}
        \dip{Teorema}{teo:bezout-equivalenze}
        \dip{Teorema}{teo:lemma-euclide}
    \end{prerequisiti}
    Siano $a, b \in \mbZ \setminus \{0\}$, $d \in \mathrm{MCD}(a, b)$, e sia $(x_0, y_0)$ una soluzione di $ax + by = c$. Allora l'insieme delle soluzioni è
    \[ \Big\{\Big(x_0 + \frac{b}{d}k,\ y_0 - \frac{a}{d}k\Big) \ \Big|\ k \in \mbZ\Big\}. \]
\end{Teorema}
\begin{dimostrazione}
    \begin{caso}{Sono soluzioni}
        $a\big(x_0 + \frac{b}{d}k\big) + b\big(y_0 - \frac{a}{d}k\big) = ax_0 + by_0 + \frac{ab}{d}k - \frac{ab}{d}k = c$.
    \end{caso}
    \begin{caso}{Sono tutte}
        Sia $(x, y)$ una soluzione: $ax + by = c = ax_0 + by_0$. Dividendo per $d$ e riordinando: $\frac{a}{d}(x - x_0) = \frac{b}{d}(y_0 - y)$. Gli interi
        $\frac{a}{d}$ e $\frac{b}{d}$ sono coprimi (come $a_1$ e $b_1$ nel teorema precedente), quindi per il lemma di Euclide $\frac{b}{d} \mid x - x_0$ e
        $\frac{a}{d} \mid y_0 - y$: $x - x_0 = k\frac{b}{d}$ e $y_0 - y = h\frac{a}{d}$. Sostituendo, $\frac{a}{d}\frac{b}{d}k = \frac{b}{d}\frac{a}{d}h$, quindi $h = k$, e
        $(x, y) = (x_0 + \frac{b}{d}k, y_0 - \frac{a}{d}k)$.
    \end{caso}
\end{dimostrazione}
\end{concetto}

\begin{spiegazione}
    Le soluzioni sono i \textbf{punti a coordinate intere} sulla retta $ax + by = c$, e si susseguono a passi regolari di
    $\big(\frac{b}{d}, -\frac{a}{d}\big)$. Per $2x + 3y = 7$ ($d = 1$), partendo dalla soluzione $(2, 1)$ (infatti $4 + 3 = 7$), le soluzioni sono
    $(2 + 3k, 1 - 2k)$:
    \begin{center}
    \begin{tikzpicture}[scale=0.42]
        \draw[step=1, black!10] (-5,-4) grid (9,6);
        \draw[-{Stealth}, black!50] (-5,0) -- (9.5,0) node[right, font=\scriptsize] {$x$};
        \draw[-{Stealth}, black!50] (0,-4) -- (0,6.5) node[above, font=\scriptsize] {$y$};
        \draw[primary!70!black, line width=0.8pt] (-5,17/3) -- (9,-11/3);
        \foreach \x/\y in {-4/5, -1/3, 2/1, 5/-1, 8/-3} \fill[orange!80!black] (\x,\y) circle (5pt);
        \draw[-{Stealth[length=2mm]}, seagreen!65!black, line width=1.1pt] (2,1) -- (5,-1);
        \node[font=\scriptsize, seagreen!60!black, anchor=south west] at (3.5,0.05) {$(+3, -2)$};
        \node[font=\scriptsize, anchor=west] at (2.2,1.5) {$(2, 1)$};
        \node[font=\small, align=left, anchor=west] at (10.5,1)
            {la retta $2x + 3y = 7$ passa per\\ infiniti punti, ma solo alcuni\\ hanno coordinate intere:\\
             $\dots, (-1, 3), (2, 1), (5, -1), (8, -3), \dots$};
    \end{tikzpicture}
    \end{center}
    Il passo non può essere più corto: per restare sulla retta, se $x$ aumenta di $t$, $3y$ deve diminuire di $2t$, quindi $3 \mid 2t$ e
    (lemma di Euclide) $3 \mid t$. Invece $2x + 4y = 5$ non ha soluzioni ($2 \nmid 5$): la retta non passa per nessun punto intero.
\end{spiegazione}

% ──────────────────────────────────────────────────────────────
\section{Equazioni congruenziali}\label{sec:congruenziali}

\dfn[congruenziale]{Equazione congruenziale di primo grado}{
    Siano $m \in \mbZ \setminus \{0\}$ e $a, b \in \mbZ$. La funzione $[n]_m \in \mbZ_m \mapsto [an - b]_m \in \mbZ_m$ si chiama \textbf{equazione
    congruenziale} di primo grado in un'incognita, di termini $a, b$ e modulo $m$; si scrive $ax \equiv_m b$. Un intero $n$ si dice
    \textbf{soluzione} se $an \equiv_m b$, cioè se $[a]_m [n]_m = [b]_m$.
}

Se $n$ è una soluzione, lo sono tutti gli elementi di $[n]_m$ (perché $[a][n]$ dipende solo dalla classe di $n$): le soluzioni si danno
come classi di resto. Ad esempio $4x \equiv_2 3$ non ha soluzioni ($4n$ è pari, $3$ è dispari).

\begin{concetto}
\begin{Teorema}{Quando un'equazione congruenziale ha soluzioni}{congruenziale-soluzioni}
    \begin{prerequisiti}
        \dip{Definizione}{def:congruenziale}
        \dip{Teorema}{teo:bezout-equivalenze}
        \dip{Proposizione}{prop:mcd-associati}
    \end{prerequisiti}
    Siano $a, b \in \mbZ$, $m \in \mbZ \setminus \{0\}$ e $d \in \mathrm{MCD}(a, m)$. L'equazione $ax \equiv_m b$ ha soluzioni se e solo se $d \mid b$.
\end{Teorema}
\begin{dimostrazione}
    $n$ è soluzione se e solo se $m \mid b - an$, cioè se e solo se esiste $y$ con $b - an = my$, cioè $an + m(-y) = b$. Quindi
    $ax \equiv_m b$ ha soluzioni se e solo se l'equazione diofantea $ax + my = b$ ne ha, e per il Teorema~\ref{teo:bezout-equivalenze}, punto 4,
    questo accade se e solo se $d \mid b$. (Se $a = 0$, $\mathrm{MCD}(0, m) = \{\pm m\}$ e l'equazione $0 \equiv_m b$ ha soluzioni se e solo se $m \mid b$.)
\end{dimostrazione}
\end{concetto}

\begin{spiegazione}
    Un'equazione congruenziale \textbf{è} un'equazione diofantea travestita: «$an$ e $b$ differiscono per un multiplo di $m$» vuol dire
    «$an + my = b$ per qualche intero $y$». Ad esempio $12x \equiv_9 3$ ha soluzioni ($\mathrm{MCD}(12, 9) \ni 3$ e $3 \mid 3$), mentre $12x \equiv_9 5$ no
    ($3 \nmid 5$): $12n$ è sempre multiplo di $3$, e lo è anche $12n - 9y$, che non può valere $5$.
\end{spiegazione}

\begin{concetto}
\begin{Corollario}{Gli invertibili di $\mbZ_m$}{invertibili-zm}
    \begin{prerequisiti}
        \dip{Teorema}{teo:congruenziale-soluzioni}
        \dip{Definizione}{def:coprimi}
        \dip{Teorema}{teo:divisori-cancellabili}
        \dip{Teorema}{teo:invertibile-cancellabile}
    \end{prerequisiti}
    Sia $m \in \mbZ \setminus \{0\}$ e $a \in \mbZ$. Sono equivalenti:
    \begin{enumerate}
        \item $[a]_m$ è invertibile in $(\mbZ_m, \cdot)$;
        \item $a$ e $m$ sono coprimi;
        \item $[a]_m$ non è un divisore dello zero in $(\mbZ_m, +, \cdot)$.
    \end{enumerate}
\end{Corollario}
\begin{dimostrazione}
    \begin{caso}{(1) $\leftrightarrow$ (2)}
        $[a]_m$ è invertibile se e solo se l'equazione $ax \equiv_m 1$ ha soluzioni, cioè (Teorema~\ref{teo:congruenziale-soluzioni}) se e solo se
        $d \mid 1$ per $d \in \mathrm{MCD}(a, m)$, cioè se e solo se $1 \in \mathrm{MCD}(a, m)$.
    \end{caso}
    \begin{caso}{(1) $\rightarrow$ (3)}
        Se $[a][b] = [0]$ con $[a]$ invertibile, moltiplicando per l'inverso si ottiene $[b] = [0]$: non esiste $[b] \neq [0]$ con $[a][b] = [0]$.
    \end{caso}
    \begin{caso}{(3) $\rightarrow$ (2), per contrapposizione}
        Se $a$ e $m$ non sono coprimi, esiste $d \in \mathrm{MCD}(a, m)$ con $d > 1$ (possiamo prenderlo positivo). Scriviamo $a = a'd$. Allora
        $[\frac{m}{d}] \neq [0]$ (perché $0 < |\frac{m}{d}| < |m|$) e $[a][\frac{m}{d}] = [a'd\frac{m}{d}] = [a'm] = [0]$: $[a]$ è un divisore dello zero.
    \end{caso}
\end{dimostrazione}
\end{concetto}

\begin{spiegazione}
    In $\mbZ_m$ ogni classe non nulla è \textbf{o invertibile o divisore dello zero}, e lo si decide guardando il massimo comun divisore
    con $m$. Sull'orologio di $\mbZ_{12}$ (in verde gli invertibili, in rosso i divisori dello zero diversi da $[0]$):
    \begin{center}
    \begin{tikzpicture}
        \draw[black!40, line width=0.8pt] (0,0) circle (1.45);
        \foreach \r in {0,...,11} {
            \pgfmathsetmacro{\ang}{90 - 30*\r}
            \ifnum\r=0 \def\col{black!8}\else
                \ifnum\r=1 \def\col{seagreen!35}\else \ifnum\r=5 \def\col{seagreen!35}\else \ifnum\r=7 \def\col{seagreen!35}\else
                \ifnum\r=11 \def\col{seagreen!35}\else \def\col{cherryred!22}\fi\fi\fi\fi\fi
            \node[circle, fill=\col, draw=black!40, minimum size=7mm, inner sep=0pt, font=\small] at (\ang:1.45) {$\r$};
        }
        \node[font=\small, align=left, anchor=west] at (2.3,0.3)
            {invertibili: $1, 5, 7, 11$ (coprimi con $12$),\\ ad esempio $[5][5] = [25] = [1]$;\\[1mm]
             divisori dello zero: $2, 3, 4, 6, 8, 9, 10$,\\ ad esempio $[8][3] = [24] = [0]$ con $[3] \neq [0]$.};
    \end{tikzpicture}
    \end{center}
    Per $m$ primo tutte le classi non nulle sono coprime con $m$: ritroviamo che $\mbZ_p$ è un campo.
\end{spiegazione}

% ──────────────────────────────────────────────────────────────
\section{Risolvere un'equazione congruenziale}\label{sec:risolvere-congruenziali}

\begin{concetto}
\begin{Proposizione}{Operazioni che non cambiano le soluzioni}{semplificazioni}
    \begin{prerequisiti}
        \dip{Definizione}{def:congruenziale}
        \dip{Proposizione}{prop:congruenza-z}
        \dip{Corollario}{cor:invertibili-zm}
    \end{prerequisiti}
    Sia $e$ l'equazione $ax \equiv_m b$. Le equazioni seguenti hanno le stesse soluzioni (intere) di $e$:
    \begin{enumerate}
        \item $a'x \equiv_m b'$, se $a' \equiv_m a$ e $b' \equiv_m b$ (ad esempio $a' = a \bmod m$, $b' = b \bmod m$);
        \item $(ak)x \equiv_{mk} bk$, per ogni $k \in \mbZ \setminus \{0\}$; e quindi anche, se $a = a'k$, $b = b'k$, $m = m'k$, l'equazione $a'x \equiv_{m'} b'$;
        \item $(al)x \equiv_m bl$, per ogni $l$ coprimo con $m$.
    \end{enumerate}
\end{Proposizione}
\begin{dimostrazione}
    \begin{caso}{Punto 1}
        $an \equiv_m a'n$ e $b \equiv_m b'$, quindi $an \equiv_m b \leftrightarrow a'n \equiv_m b'$ (transitività e simmetria).
    \end{caso}
    \begin{caso}{Punto 2}
        $mk \mid bk - akn \leftrightarrow mk \mid (b - an)k \leftrightarrow m \mid b - an$, perché $k \neq 0$ si può cancellare. La seconda parte è la prima
        letta al contrario, con $a', b', m'$ al posto di $a, b, m$.
    \end{caso}
    \begin{caso}{Punto 3}
        $[l]_m$ è invertibile (Corollario~\ref{cor:invertibili-zm}); moltiplicando $[a][n] = [b]$ per $[l]$ o per $[l]^{-1}$ si passa da una forma all'altra:
        $[a][n] = [b] \leftrightarrow [al][n] = [bl]$.
    \end{caso}
\end{dimostrazione}
\end{concetto}

\begin{spiegazione}
    Sono le stesse mosse che si fanno con le equazioni «normali» (semplificare, moltiplicare per un numero non nullo), con una
    precisazione importante: si può moltiplicare o dividere \textbf{solo i due membri} per un numero \textbf{invertibile modulo $m$}
    (punto 3), mentre per dividere per un fattore comune non invertibile bisogna dividere \textbf{anche il modulo} (punto 2). Ad
    esempio $2x \equiv_6 4$ e $x \equiv_3 2$ hanno le stesse soluzioni ($\dots, -1, 2, 5, 8, \dots$), ma $x \equiv_6 2$ no: le perde la metà.
\end{spiegazione}

\begin{concetto}
\begin{Teorema}{Metodo risolutivo}{metodo-congruenziali}
    \begin{prerequisiti}
        \dip{Proposizione}{prop:semplificazioni}
        \dip{Teorema}{teo:congruenziale-soluzioni}
        \dip{Corollario}{cor:invertibili-zm}
        \dip{Teorema}{teo:bezout}
    \end{prerequisiti}
    Siano $m \neq 0$, $a, b \in \mbZ$, $d \in \mathrm{MCD}(a, m)$ con $d \mid b$. Scriviamo $a = a'd$, $b = b'd$, $m = m'd$ e sia $l$ un intero con
    $a'l \equiv_{m'} 1$ (esiste perché $a'$ e $m'$ sono coprimi). Allora le soluzioni di $ax \equiv_m b$ sono esattamente gli interi della classe
    $[b'l]_{m'}$: una classe di resto modulo $\frac{m}{d}$.
\end{Teorema}
\begin{dimostrazione}
    Per la Proposizione~\ref{prop:semplificazioni}, punto 2, l'equazione ha le stesse soluzioni di $a'x \equiv_{m'} b'$. $a'$ e $m'$ sono coprimi
    (come nella dimostrazione del Teorema~\ref{teo:bezout-equivalenze}), quindi $[a']_{m'}$ è invertibile con inverso $[l]_{m'}$, che si trova con
    l'algoritmo delle divisioni successive esteso. Moltiplicando per $l$ (punto 3), le soluzioni sono quelle di $x \equiv_{m'} b'l$, cioè
    gli elementi di $[b'l]_{m'}$.
\end{dimostrazione}
\end{concetto}

\begin{spiegazione}
    Il metodo, passo per passo:
    \begin{passi}
        \item ridurre $a$ e $b$ tra $0$ e $|m| - 1$;
        \item calcolare $d \in \mathrm{MCD}(a, m)$: se $d \nmid b$ non ci sono soluzioni;
        \item dividere $a$, $b$ e $m$ per $d$;
        \item trovare l'inverso $l$ di $a'$ modulo $m'$ (algoritmo esteso, o per tentativi se $m'$ è piccolo);
        \item la soluzione è la classe $[b'l]_{m'}$.
    \end{passi}
    \textbf{Esempio 1}: $12x \equiv_7 3$. Riduco: $5x \equiv_7 3$. $\mathrm{MCD}(5, 7) = \{\pm 1\}$. L'inverso di $5$ modulo $7$ è $3$ ($5 \cdot 3 = 15 = 2 \cdot 7 + 1$).
    Soluzione: $x \equiv_7 3 \cdot 3 = 9 \equiv_7 2$. Verifica: $12 \cdot 2 = 24 = 3 \cdot 7 + 3$.

    \textbf{Esempio 2}: $12x \equiv_9 3$. Riduco: $3x \equiv_9 3$. $d = 3 \mid 3$; divido per $3$: $x \equiv_3 1$. Soluzioni: $[1]_3 = \{\dots, -2, 1, 4, 7, 10, \dots\}$,
    cioè tre classi modulo $9$: $[1]_9$, $[4]_9$, $[7]_9$.

    \textbf{Esempio 3}: $100x \equiv_{11} 2$. Poiché $100 = 9 \cdot 11 + 1$, è $x \equiv_{11} 2$.
\end{spiegazione}

% ──────────────────────────────────────────────────────────────
\section{Periodo di un elemento di un gruppo}\label{sec:periodo}

\dfn[periodo]{Periodo di un elemento}{
    Sia $(g, *)$ un gruppo con neutro $1_g$ e sia $x \in g$. Se esiste $n \geq 1$ con $x^n = 1_g$, il più piccolo di tali $n$ si chiama
    \textbf{periodo} (o \textbf{ordine}) di $x$; altrimenti si dice che $x$ ha periodo infinito. In notazione additiva il periodo è il più
    piccolo $n \geq 1$ con $nx = 0$.
}

\begin{concetto}
\begin{Teorema}{Quando due potenze coincidono}{potenze-coincidono}
    \begin{prerequisiti}
        \dip{Definizione}{def:periodo}
        \dip{Definizione}{def:potenze}
        \dip{Teorema}{teo:divisione-euclidea}
        \dip{Proposizione}{prop:cicliche}
    \end{prerequisiti}
    Sia $(g, *)$ un gruppo e sia $x \in g$.
    \begin{enumerate}
        \item Se $x$ ha periodo $n$, allora per ogni $h, k \in \mbZ$: \ $x^h = x^k \leftrightarrow h \equiv_n k$. In particolare
              $\langle x \rangle = \{1_g, x, x^2, \dots, x^{n - 1}\}$ ha esattamente $n$ elementi.
        \item Se $x$ ha periodo infinito, le potenze $x^h$ ($h \in \mbZ$) sono a due a due distinte.
    \end{enumerate}
\end{Teorema}
\begin{dimostrazione}
    Per le regole delle potenze con esponente intero, $x^h = x^k \leftrightarrow x^{h - k} = 1_g$.
    \begin{caso}{Punto 1}
        Dividiamo: $h - k = nq + r$ con $0 \leq r < n$. Allora $x^{h - k} = (x^n)^q * x^r = 1_g^q * x^r = x^r$. Quindi $x^{h - k} = 1_g \leftrightarrow x^r = 1_g$, e per
        la minimalità di $n$ (con $0 \leq r < n$) questo accade se e solo se $r = 0$, cioè $n \mid h - k$. Per la Proposizione~\ref{prop:cicliche}
        $\langle x \rangle = \{x^h \mid h \in \mbZ\}$; ogni $x^h$ è uguale a $x^{h \bmod n}$, e $x^0, \dots, x^{n - 1}$ sono distinti (i loro esponenti non sono congrui modulo $n$).
    \end{caso}
    \begin{caso}{Punto 2}
        Se fosse $x^h = x^k$ con $h \neq k$, si avrebbe $x^{|h - k|} = 1_g$ con $|h - k| \geq 1$ (se $x^{h - k} = 1_g$ anche il suo inverso $x^{k - h}$ vale $1_g$):
        $x$ avrebbe periodo finito.
    \end{caso}
\end{dimostrazione}
\end{concetto}

\begin{spiegazione}
    Le potenze di un elemento di periodo $n$ girano in \textbf{cerchio}: dopo $n$ passi si torna a $1_g$ e tutto si ripete, quindi
    l'esponente conta solo modulo $n$. Per un elemento di periodo $6$:
    \begin{center}
    \begin{tikzpicture}
        \foreach \i in {0,...,5} {
            \pgfmathsetmacro{\ang}{90 - 60*\i}
            \node[circle, fill=primary!8, draw=primary!60!black, minimum size=8mm, inner sep=0pt, font=\small] (p\i) at (\ang:1.4) {$x^{\i}$};
        }
        \foreach \i [evaluate=\i as \j using {int(mod(\i+1,6))}] in {0,...,5}
            \draw[-{Stealth[length=2mm]}, orange!80!black, line width=0.7pt] (p\i) to[bend left=18] (p\j);
        \node[font=\small, align=left, anchor=west] at (2.3,0.2)
            {ogni freccia moltiplica per $x$;\\ $x^6 = x^0 = 1_g$, \ $x^7 = x^1$, \ $x^{20} = x^{2}$ ($20 \equiv_6 2$),\\
             $x^{-1} = x^5$ ($-1 \equiv_6 5$).};
    \end{tikzpicture}
    \end{center}
    Ad esempio in $(\mbZ_{12}, +)$ la classe $[4]$ ha periodo $3$: $[4], [8], [12] = [0]$. E se $g = \langle x \rangle$ ha $19$ elementi, $x$ ha periodo $19$ e
    $(x^5)^k = x$ se e solo se $5k \equiv_{19} 1$, cioè $k \equiv_{19} 4$ ($5 \cdot 4 = 20$): $(x^5)^4 = x^{20} = x$. Quindi anche $x^5$ genera tutto $g$.
\end{spiegazione}
